% ============================================================
% ch01.tex —— 第 1 章 余数王国：除法竖式里被扔掉的宝贝
% ============================================================
\chapter{余数王国：除法竖式里被扔掉的宝贝}

你做过几百道除法竖式了吧？每一道竖式的结尾，你都写下两个数：商和余数。商永远是主角，余数是边角料——老师甚至叮嘱你"除不尽要验算余数，别忘了加回去"。从今天起，立场反转：\textbf{本章的主角是余数}。我们将搬进一个由余数组成的小王国，你会发现，加法、乘法在这个小世界里照常营业，而一些吓人的大数问题在这里变得像掰手指头一样简单。

\begin{kaiwei}
今天是星期六。$100$ 天后是星期几？$10^{100}$ 天后呢？
\end{kaiwei}

第一问你在序章已经会了：$100 = 14\times 7 + 2$，余数 $2$ 说了算，星期一。第二问好像要人命：$10^{100}$ 是 $1$ 后面跟 $100$ 个零的天文数字，宇宙的年龄才 $4.4 \times 10^{17}$ 秒左右——这个天数连"天文学数字"都不配形容。但等一等——我们关心的根本不是这个数有多大，而是\textbf{它除以 $7$ 余几}！$10^{100}$ 天后的星期几只由这个余数决定。怎么求？硬除显然不可能，本章结束时你会自己算出来，全程不超过五步。

\section{把整数分房间}

先立规矩。做除法 $a \div b$（$b > 0$），我们要求商 $q$ 和余数 $r$ 满足
\[
a = bq + r, \qquad 0 \le r < b .
\]
这叫\Keyword{带余除法}。规矩的重点在右边那个不等式：\textbf{余数必须比除数小}，而且不小于 $0$。这两个要求合起来保证了：对固定的除数，\textbf{每个整数的余数是唯一的}——你不可能算出"$23$ 除以 $7$ 余 $2$"又算出"余 $9$"，因为 $9$ 比 $7$ 大，还能再商一次。

于是，除以 $3$，余数只能是 $0, 1, 2$ 三种；除以 $7$，余数只有 $0$ 到 $6$ 七种。这意味着一件奇妙的事：\textbf{按"除以 $3$ 的余数"分房间，全体整数——包括负数，包括那些大得没边的数——只住三个房间}：

\begin{center}
\begin{tabular}{cll}
\toprule
房间 & 住户 & 特征 \\
\midrule
$0$ 号房 & $\cdots, -6, -3, 0, 3, 6, 9, 12, \cdots$ & 3 的倍数 \\
$1$ 号房 & $\cdots, -5, -2, 1, 4, 7, 10, 13, \cdots$ & 除以 3 余 1 \\
$2$ 号房 & $\cdots, -4, -1, 2, 5, 8, 11, 14, \cdots$ & 除以 3 余 2 \\
\bottomrule
\end{tabular}
\end{center}

\textbf{负数也乖乖住进来了}，而且规矩不破。看 $-5$：它除以 $3$ 余几？按老习惯写 $-5 = 3 \times (-2) + 1$——商取 $-2$，余数 $1$，落在 $0 \le 1 < 3$ 内。所以 $-5$ 住 $1$ 号房。注意\textbf{不能}写 $-5 = 3\times(-1) - 2$：那个"$-2$"是负的，违反"余数非负"的规矩。再练两个：$-18 = 3\times(-6) + 0$，住 $0$ 号房；$-1 = 3\times(-1) + 2$，住 $2$ 号房。房间编号像钟面一样循环，没有"最外面"。

\begin{faming}
两个数如果除以 $n$ 余数相同，就说它们\Keyword{在模 $n$ 意义下相等}（土名：住同一间房），记作
\[
a \equiv b \pmod{n}.
\]
比如 $13 \equiv 6 \pmod 7$（都余 $6$），$-5 \equiv 1 \pmod 3$（都余 $1$），$10^{100} \equiv ? \pmod 7$（待会儿揭晓）。正式名字叫\textbf{同余}，是数学家高斯 24 岁那年（1801 年出版《算术研究》）发明的记号——那本书把数论从零散的谜题集变成了系统学科，而你在这一页已经把它的第一个记号用明白了。
\end{faming}

\section{房间的加法和乘法：余数替换原理}

住同一间房的数，可以互相替换吗？比如 $100 \equiv 2 \pmod 7$，那"$100$ 天后"和"$2$ 天后"的星期几确实一样——生活经验支持替换。但这不是巧合，而是可以证明的定理。它如此重要，值得起个名字：

\begin{dingli}[余数替换原理]
若 $a \equiv b \pmod n$ 且 $c \equiv d \pmod n$，则
\[
a + c \equiv b + d \pmod n, \qquad a \times c \equiv b \times d \pmod n .
\]
\end{dingli}

\begin{proof}
一句话：\textbf{房间里怎么加减乘，出来再验余数，结果不会错}。写清楚些。"$a$ 与 $b$ 同房"意味着 $a = b + nk$（$k$ 是某个整数，多出来的整份）；"$c$ 与 $d$ 同房"意味着 $c = d + nm$。于是

加法：
\[
a + c = (b + nk) + (d + nm) = (b + d) + n(k + m),
\]
多出来的两份 "$n$ 的倍数" 打包成一份，所以 $a + c$ 与 $b+d$ 同房。

乘法：
\[
a \times c = (b + nk)(d + nm) = bd + \underline{bnm} + \underline{nkd} + \underline{n^2 km},
\]
三项下划线的部分，每一项都含因子 $n$（第一项有 $n$，第二项有 $n$，第三项有 $n^2 = n\cdot n$），所以
\[
a\times c = bd + n(bm + kd + nkm),
\]
括号里是个整数，于是 $a\times c$ 与 $bd$ 同房。证毕。
\end{proof}

这条原理看似平淡，威力却大。它的本质是：\textbf{在"问余数"的问题里，你可以随时把任何数换成同房的数（通常是更小的），绝不影响答案}。所有接下来的戏法，都是这一句话的变奏。

先来个小的练手：求 $2^{10}$ 除以 $7$ 的余数。硬算是可以的——$2^{10} = 1024$，$1024 = 146\times 7 + 2$，余 $2$。但聪明人不硬算，而是\textbf{边乘边压扁}：
\[
2^1 \equiv 2,\qquad 2^2 = 4 \equiv 4,\qquad 2^3 = 8 \equiv 1 \pmod 7 .
\]
$2^3$ 的余数已经\textbf{转回了 $1$}！好消息：一旦出现 $1$，后面就是循环——$2^4 = 2^3\times 2 \equiv 1\times 2 = 2$，$2^5 \equiv 4$，$2^6 \equiv 1$……周期是 $3$：$2, 4, 1, 2, 4, 1, \cdots$。于是
\[
2^{10} = 2^{9}\times 2 = (2^3)^3 \times 2 \equiv 1^3 \times 2 = 2 \pmod 7 .
\]
和硬算的对上了，但注意中间步骤里\textbf{从未出现超过 $7$ 的数}——哪怕指数是一亿，这台机器也照转不误。

\textbf{余数的世界里有周期。}这是本章最重要的发现。为什么会有周期？因为世界太小：模 $7$ 的非零余数只有 $6$ 个房间，$2$ 的幂每乘一次搬一次家，房间有限，\textbf{迟早撞回住过的房间}；而一旦 $2^i \equiv 2^j$（$i<j$），由替换原理两边约去 $2^i$（这个"约去"的合法性要等到第 10 章才完全清白——模 $7$ 是素数的世界，第 6 章会看到它为什么特殊），就得到 $2^{j-i} \equiv 1$：撞房意味着出现 $1$，出现 $1$ 意味着开始循环。

\begin{cidian}
土名"余数的世界""压扁"$\to$ 正式名\textbf{模 $n$ 算术}（modular arithmetic）；"取模"就是"取余数"；符号 $\equiv$ 读"同余于"，$\pmod n$ 读"模 $n$"。高斯的原文用三条横线，意思是"这是一种加强版的相等"——相等的数处处相等，同余的数在余数世界里处处等价。
\end{cidian}

\section{个位数字循环：一场小型余数表演}

\begin{kaiwei}
$2$ 的幂的个位数字依次是 $2, 4, 8, 6, 2, 4, 8, 6, \cdots$，四个数字打转。这个循环会永远持续吗？$2^{100}$ 的个位数字是多少？
\end{kaiwei}

第一步是想通一件事：\textbf{个位数字是什么？}它就是"除以 $10$ 余几"——个位是 $7$ 的数，除以 $10$ 余 $7$；个位是 $0$，余 $0$。所以"个位数字问题"恰好是"模 $10$ 的余数问题"！替换原理立刻上岗：$2^n$ 的个位由 $2^n \bmod 10$ 决定。手算一遍，边乘边压：
\[
\begin{array}{cccccccccc}
2^1 & 2^2 & 2^3 & 2^4 & 2^5 & 2^6 & 2^7 & 2^8 & \cdots \\
2 & 4 & 8 & 16\to 6 & 32\to 2 & 64\to 4 & 128\to 8 & 256\to 6 & \cdots
\end{array}
\]
周期是 $4$：$2, 4, 8, 6$。为什么一定会有周期？模 $10$ 的世界只有 $10$ 个房间，幂迟早撞房——同一个论证，第二次生效。

$2^{100}$ 的个位：$100 = 4\times 25$，恰好在周期的第 $4$ 个位置上（周期起点是 $2^1$，所以看余数：$100$ 除以 $4$ 余 $0$，对应周期末位），答案是 $\boxed{6}$。这里"余 $0$ 对应最后一个"是新手最容易栽的坑：$2^4, 2^8, 2^{12}, \cdots, 2^{100}$ 的个位全是 $6$，而 $2^1, 2^5, 2^9$（指数除以 $4$ 余 $1$）的个位全是 $2$。\textbf{用指数除以周期的余数查表}，指数再大也只是一次除法。

同样的戏法对别的底数照演，把表格补全（这是本章最值得抄进笔记本的一张表）：

\begin{center}
\begin{tabular}{ccc}
\toprule
底数 & 个位循环 & 周期 \\
\midrule
$2$ & $2,4,8,6$ & $4$ \\
$3$ & $3,9,7,1$ & $4$ \\
$4$ & $4,6$ & $2$ \\
$5$ & $5$ & $1$ \\
$6$ & $6$ & $1$ \\
$7$ & $7,9,3,1$ & $4$ \\
$8$ & $8,4,2,6$ & $4$ \\
$9$ & $9,1$ & $2$ \\
\bottomrule
\end{tabular}
\end{center}

看表找乐子：$5$ 和 $6$ 的个位纹丝不动（周期 $1$）；$4$ 和 $9$ 只玩二人转；$2,3,7,8$ 整整齐齐都是周期 $4$。这些"为什么"要到第 6、9 章才完全揭开（周期其实整除 $4$，而 $4$ 恰好是"$10$ 的世界"容许的最大周期）——但你现在已经能\textbf{用}这张表秒杀一类题：$7^{99}$ 的个位？$99 = 4\times 24 + 3$，余 $3$，查 $7$ 的循环第 $3$ 位：$3$。

\section{兑现承诺：$10^{100}$ 天后是星期几？}

工具齐了，现在收掉开胃问题的尾巴。目标是 $10^{100} \bmod 7$。策略：先找 $10$ 的幂在模 $7$ 世界里的周期。

\textbf{第一步}，把 $10$ 压扁：$10 = 7 + 3$，所以
\[
10 \equiv 3 \pmod 7 .
\]

\textbf{第二步}，逐个算幂（边乘边压，随时把大数换成同房的小数）：
\[
10^2 \equiv 3^2 = 9 \equiv 2; \qquad
10^3 \equiv 2 \times 3 = 6 \equiv -1 \pmod 7 .
\]
这里出了个宝贝：$10^3 \equiv -1$。负一同余是个好东西，因为\textbf{平方必得正一}：

\textbf{第三步}，
\[
10^6 \equiv (-1)^2 = 1 \pmod 7 .
\]
周期找到了：$6$。（它确实不可能更短——你可以自己验证 $10^1, 10^2, 10^3$ 都不是 $1$。）

\textbf{第四步}，用周期砍指数。$100 = 6\times 16 + 4$，所以
\[
10^{100} = \left(10^{6}\right)^{16} \times 10^{4} \equiv 1^{16} \times 10^{4} \pmod 7 .
\]

\textbf{第五步}，只剩小尾巴 $10^4$：
\[
10^4 = 10^3 \times 10 \equiv (-1)\times 3 = -3 \equiv 4 \pmod 7 .
\]

结论：$10^{100} \equiv 4 \pmod 7$。星期六再过 $4$ 天——星期日、一、二、三——是\textbf{星期三}。一个 $101$ 位数的问题，五步之内解决，中间出现的最大的数是 $100$。\textbf{这就是搬进余数王国的房租回报：数的巨无霸，在余数世界里只有身份证号，没有体重。}

\begin{dongshou}
\begin{itemize}
  \yisi 把下列负数安置进"除以 $5$"的房间：$-18$、$-3$、$-25$。（提示：允许商取负数，例 $-18 = 5\times(-4) + 2$，住 $2$ 号房。）
  \yier 求 $3^{100}$ 的个位数字；再求 $7^{99}$、$8^{2026}$ 的个位数字。（全部查上面的周期表，先算指数除以周期的余数。$8$ 的周期是 $4$：$2026 = 4\times 506 + 2$，查第 $2$ 位。）
  \yier 求 $2^{2026}$ 除以 $7$ 的余数。（先确认 $2^3 \equiv 1 \pmod 7$，再用 $2026 = 3\times 675 + 1$。答案：$2$。）
  \yier 今天星期三，$2^{100}$ 天后是星期几？（把本章两台机器串起来用：先求 $2^{100} \bmod 7$——注意 $2^3 \equiv 1$ 的周期是 $3$ 不是 $6$——再从星期三往后数。）
\end{itemize}
\end{dongshou}

\begin{shaonao}
在余数王国（模 $10$）里，$2 \times 5 = 10 \equiv 0$：两个\textbf{都不是零}的数，乘起来竟然"归零"了！这破坏了我们对乘法的一条朴素信仰。请你系统地侦察：在模 $6$ 的世界里，找出所有"两个非零数乘出零"的组合；再试模 $8$、模 $12$。然后回答两个问题：(1) 这些"惹祸"的数有什么共同点？（和模数共享因数。）(2) 模 $7$ 的世界里，这样的事故一次都发不出来——为什么偏偏是 $7$？记住你的发现，第 10 章它会变成一颗重磅炸弹：\textbf{有这种事故的世界，除法注定瘫痪}。
\end{shaonao}

\begin{shaonao}
减法在余数王国也成立：若 $a\equiv b \pmod n$，则对任何 $c$，$a - c \equiv b - c \pmod n$。试着证明它（把 $a$ 写成 $b + nk$，两边同减 $c$ 即可）。然后用它解释一个生活现象：如果今天是星期二，那么"上周三"是多少天前？按"天数可以为负"的规则，是 $-6$ 天；而 $-6 \equiv 1 \pmod 7$，意思是"上周三"和"昨天"（星期一，$-1$ 天……等等，$-1 \equiv 6$？）——请你把这笔账彻底算清楚：$-6$ 和 $-1$ 在模 $7$ 的世界里同房吗？星期二往前推 $6$ 天和往前推 $1$ 天，差整整一个星期，日子却都是星期三和星期一\textbf{不对}……亲手排一遍日历，看看哪一步想岔了。做对这道题，你就真正习惯了"负的天数也带余数身份"。
\end{shaonao}

\begin{chengshi}
严格的教科书会把"余数替换原理"完整地证明给你，包括减法和除法（除法要小心！$6 \equiv 0 \pmod 3$，两边除以 $2$ 得 $3 \equiv 0 \pmod 3$——错！$3 \bmod 3 = 0$ 没错……再试 $8 \equiv 2 \pmod 6$，除以 $2$ 得 $4 \equiv 1 \pmod 6$，但 $4 \bmod 6 = 4 \ne 1$——\textbf{同余两边不能随便除}！）。本册只用加法与乘法的替换；"约去公因子合法的精确条件"（要求逆元）放在第 10 章。另外，"周期一定存在"我们用"房间有限迟早撞房"论证了，但"撞房就意味着回到 $1$"这一步，在模数不是素数时需要小心（撞房只保证回到某个循环，不一定从 $1$ 开始）——第 9 章的拉格朗日定理会把这件事一次说透。
\end{chengshi}

你已经学会了把数压扁。下一章要解决一个更古老的问题：两个数，怎么找到它们"最大的共同尺寸"？别被这句话绕晕——直接上工地。
